\section{Bridge L2 Senza Gateway}\label{bridge-l2-senza-gateway}

Questa architettura prevede solo la creazione di un bridge L2, per
istanziare una rete LAN virtuale sui cui attaccare i container da
verificare. I container duplicati mantengono la topologia di rete
invariata, e possono comunicare tra di loro, tuttavia il bridge L2 non
ha un gateway verso cui instradare per accedere alla WAN. Dunque al
livello 3 i container non possono essere visti dall'esterno della LAN.

Questa architettura è molto semplice da implementare ed è anche
efficiente, in quanto fa uso solo di strutture dati primitive del kernel
linux e di un unico bridge.

Tuttavia rende impossibile per i container l'instradamento verso
l'esterno della rete. Dunque, dall'interno dei container di verifica,
comandi come \passthrough{\lstinline!apt!}, e
\passthrough{\lstinline!ping!} verso host esterni falliranno.

Questa architettura è perfetta per vulnerabilità e script di verifica
che non hanno bisogno di accedere ad internet. Immaginando però che per
qualsiasi motivo, internet potesse essere fondamentale, anche in fase di
verifica, si è scelto di scartare questa opzione.
